% Generated draft: requires mathematical and source review.
\documentclass[11pt]{article}
\usepackage[margin=25mm]{geometry}
\usepackage{fontspec}
\setmainfont{Latin Modern Roman}
\usepackage{amsmath,amssymb,mathtools,unicode-math}
\setmathfont{Latin Modern Math}
\usepackage{enumitem,needspace,xurl,hyperref}
\setlength{\parindent}{0pt}
\setlength{\parskip}{5pt}
\setlength{\emergencystretch}{3em}
\newcommand{\N}{\mathbb{N}}
\newcommand{\Z}{\mathbb{Z}}
\newcommand{\Q}{\mathbb{Q}}
\newcommand{\R}{\mathbb{R}}
\newcommand{\C}{\mathbb{C}}
\newcommand{\F}{\mathbb{F}}
\newcommand{\A}{\mathbb{A}}
\newcommand{\PP}{\mathbb{P}}
\newcommand{\E}{\mathbb{E}}
\newcommand{\eps}{\varepsilon}
\newcommand{\dd}{\,\mathrm{d}}
\begin{document}
\section*{UnsolvedMath ID: 30000447}
\textbf{Source problem code:} catalogue entry
\subsection{Definitions and mathematical statement}
Let $G$ be a finite group, let $p$ be a prime, and let $B$ be a $p$-block, using a splitting $p$-modular system so that ordinary and Brauer characters can be associated to the same block. Let $D$ be a defect group of $B$. Define
\[
 k(B)=|\operatorname{Irr}(B)|,\qquad
 l(B)=|\operatorname{IBr}(B)|,
\]
where $\operatorname{Irr}(B)$ consists of the ordinary irreducible characters in $B$ and $\operatorname{IBr}(B)$ of its irreducible Brauer characters, equivalently the isomorphism classes of simple modules in the block over a splitting field of characteristic $p$. Finally,
\[
 k(D)=|\operatorname{Irr}(D)|
      =\#\{\text{conjugacy classes of }D\}.
\]

The question is whether
\[
 \forall G\ \forall p\ \forall B\text{ a }p\text{-block of }G,
 \qquad l(B)=1\quad\Longrightarrow\quad k(B)\le k(D).
\]
Thus a block with exactly one irreducible Brauer character should have no more ordinary irreducible characters than its defect group has conjugacy classes. No abelianness or other restriction on $D$ is stated. The right-hand side is $k(D)$, not $|D|$; replacing it by the defect-group order would change the requested bound. The standard splitting-system terminology is explanatory and should be checked against the report's conventions.
\subsection{Short English statement}
Does a block with one irreducible Brauer character have at most as many ordinary irreducible characters as its defect group has conjugacy classes? The target is the exact $k(D)$ bound, including nonabelian defect groups.
\subsection{Sources}
\begin{itemize}
\item Local UnsolvedMath registry, permanent ID 30000447, source code catalogue entry, ``Character Bounds for Blocks with One Brauer Character.'' Proposed year 2006; proposer unspecified.
\item \emph{Representations of Finite Groups} (2007), Oberwolfach Reports, pp.\ 895--978, DOI 10.4171/owr/2006/15; \url{https://doi.org/10.4171/owr/2006/15}. The displayed report year 2007 differs from the DOI and proposed year 2006.
\item Supplied ``Brauer problems survey'': \url{https://www.imaginary.org/sites/default/files/snapshots/snapshot-2025-006.pdf}.
\item Supplied ``Cartan matrices and block conjectures'': \url{https://arxiv.org/abs/0707.3495}.
\item Statement status corrected\_verified; the supplied note says the block invariants and inequality were enclosed in balanced math delimiters.
\item Literature triage checked 2026-08-21: open. It places the $l(B)=1$ implication in the Brauer $k(B)$ landscape and reports partial results for block classes but no proof for all blocks. This contextual description must not replace the exact $k(D)$ target by the more usual defect-order bound. Created 2026-08-25; updated 2026-09-16. Current status and cited claims require independent verification.
\item Registry curation by Ulam AI / UnsolvedMath Contributors, CC BY 4.0; underlying report and literature provenance remain as supplied. Unreviewed statement-only draft.
\end{itemize}
\end{document}
